% !TEX root = ../../ctfp-print.tex

\lettrine[lhang=0.17]{圏}{の射にはさまざまな直感的な解釈が可能だが}、対象
$a$ から対象 $b$ への射が存在するならば、その二つの対象はある意味で
「関係している」という点では誰もが同意するだろう。射は、ある意味でその関係の
証明だ。これは、射が関係そのものであるポセット圏で特に明確だ。一般に、
二つの対象の間には同じ関係の「証明」が多数存在することがある。
これらの証明は、ホム集合と呼ぶ集合をなす。
対象を変化させると、対象の対から「証明」の集合への写像が得られる。
この写像は関手的だ——第一引数に関して反変、第二引数に関して共変である。
これを圏内の対象の間のグローバルな関係を確立するものとみなせる。
この関係は、ホム関手によって記述される：
\[\cat{C}(-, =) \Colon \cat{C}^\mathit{op}\times{}\cat{C} \to \Set\]
一般に、このような関手はすべて圏の対象の間の関係を確立するものとして解釈できる。
関係は二つの異なる圏 $\cat{C}$ と $\cat{D}$ を含むこともある。
そのような関係を記述する関手は以下のシグネチャを持ち、\newterm{プロ関手}（profunctor）と呼ばれる：
\[p \Colon \cat{D}^\mathit{op}\times{}\cat{C} \to \Set\]
数学者はこれを $\cat{C}$ から $\cat{D}$ へのプロ関手と言い（向きが逆であることに注意）、
スラッシュ付き矢印を記号として用いる：
\[\cat{C} \nrightarrow \cat{D}\]
プロ関手は $\cat{C}$ の対象と $\cat{D}$ の対象の間の
\newterm{証明関連関係}（proof-relevant relation）とみなせる。
集合の要素が関係の証明を表す。$p\ a\ b$ が空であるとき、$a$ と $b$ の間に関係はない。
関係は対称である必要はないことに注意しよう。

もう一つの有用な直感は、自己関手がコンテナであるというアイデアの一般化だ。
型 $p\ a\ b$ のプロ関手の値は、型 $a$ の要素をキーとする $b$ のコンテナと考えられる。
特に、ホム・プロ関手の要素は $a$ から $b$ への関数だ。

Haskell では、プロ関手は \code{dimap} というメソッドを備えた二引数型コンストラクタ
\code{p} として定義される。\code{dimap} は、最初のものが「逆方向」に進む
関数の対をリフトする：

\src{snippet01}
プロ関手の関手性から、\code{a} が \code{b} に関係するという証明があれば、
\code{c} から \code{a} への射と \code{b} から \code{d} への射が存在する限り、
\code{c} が \code{d} に関係するという証明が得られる。
あるいは、最初の関数を新しいキーを古いキーに変換するもの、
二番目の関数をコンテナの内容を変更するものと考えることもできる。

一つの圏内で作用するプロ関手については、型 $p\ a\ a$ の対角要素から
かなり多くの情報を取り出せる。$b \to a$ と $a \to c$ の射の対があれば、
$b$ が $c$ に関係することを示せる。さらに良いことに、単一の射を使って
非対角値に到達できる。例えば、射 $f \Colon a \to b$ があれば、
対 $\langle f, \idarrow[b] \rangle$ をリフトして $p\ b\ b$ から
$p\ a\ b$ へ移動できる：

\src{snippet02}
あるいは、対 $\langle \idarrow[a], f \rangle$ をリフトして
$p\ a\ a$ から $p\ a\ b$ へ移動できる：

\src{snippet03}

\section{二重自然変換}

プロ関手は関手なので、標準的な方法で自然変換を定義できる。
しかし多くの場合、二つのプロ関手の対角要素の間の写像を定義するだけで十分だ。
そのような変換は\newterm{二重自然変換}（dinatural transformation）と呼ばれる。
ただし、対角要素と非対角要素を結ぶ二つの方法を反映する可換条件を満たす必要がある。
関手圏 ${[}\cat{C}^\mathit{op}\times{}\cat{C}, \Set{]}$ の要素である
二つのプロ関手 $p$ と $q$ の間の二重自然変換は、射の族：
\[\alpha_a \Colon p\ a\ a \to q\ a\ a\]
であり、任意の $f \Colon a \to b$ に対して以下の図式が可換になる：

\begin{figure}[H]
  \centering
  \includegraphics[width=0.35\textwidth]{images/end.jpg}
\end{figure}

\noindent
これは自然性条件より真に弱いことに注意しよう。
$\alpha$ が ${[}\cat{C}^\mathit{op}\times{}\cat{C}, \Set{]}$ における自然変換であれば、
上の図式は二つの自然性の正方形と一つの関手性条件（プロ関手 $q$ が合成を保存する）
から構成できる：

\begin{figure}[H]
  \centering
  \includegraphics[width=0.4\textwidth]{images/end-1.jpg}
\end{figure}

\noindent
${[}\cat{C}^\mathit{op}\times{}\cat{C}, \Set{]}$ における自然変換 $\alpha$ の成分は、
対象の対 $\alpha_{a b}$ によって添字付けられることに注意しよう。
一方、二重自然変換は、それぞれのプロ関手の対角要素のみを写すため、
一つの対象によって添字付けられる。

\section{エンド}

これで、圏論の「代数」から「微積分」とでも言うべきものへと進む準備ができた。
エンド（およびコエンド）の微積分は、伝統的な微積分からアイデアや表記法を借用している。
特に、コエンドは無限の和または積分として理解でき、エンドは無限の積に似ている。
ディラックのデルタ関数に類似したものさえある。

エンドは極限の一般化であり、関手がプロ関手に置き換えられる。
コーン（錐）の代わりに、\newterm{くさび}（wedge）がある。
くさびの底はプロ関手 $p$ の対角要素で形成される。
くさびの頂点は対象（ここでは $\Set$ 値プロ関手を考えているので集合）であり、
辺は頂点から底の集合への関数の族だ。
この族を一つの多相関数——戻り値の型に多相な関数——として考えることができる：
\[\alpha \Colon \forall a\ .\ \mathit{apex} \to p\ a\ a\]
コーンとは異なり、くさびの中には底の頂点を結ぶ関数はない。
しかし、前に見たように、$\cat{C}$ 内の任意の射 $f \Colon a \to b$ が与えられると、
$p\ a\ a$ と $p\ b\ b$ の両方を共通集合 $p\ a\ b$ に結ぶことができる。
したがって、以下の図式が可換になることを要求する：

\begin{figure}[H]
  \centering
  \includegraphics[width=0.4\textwidth]{images/end-2.jpg}
\end{figure}

\noindent
これを\newterm{くさび条件}（wedge condition）と呼ぶ。次のように書ける：
\[p\ \idarrow[a]\ f \circ \alpha_a = p\ f\ \idarrow[b] \circ \alpha_b\]
あるいは、Haskell 記法を使うと：

\src{snippet04}
普遍的構成を進めて、$p$ のエンドを普遍くさびとして定義できる。
それは集合 $e$ と関数の族 $\pi$ からなり、頂点 $e'$ と族 $\alpha$ を持つ
他のくさびに対して、すべての三角形を可換にする一意な関数
$h \Colon e' \to e$ が存在する：
\[\pi_a \circ h = \alpha_a\]

\begin{figure}[H]
  \centering
  \includegraphics[width=0.4\textwidth]{images/end-21.jpg}
\end{figure}

\noindent
エンドの記号は積分記号で、「積分変数」が下付きの位置に来る：
\[\int_c p\ c\ c\]
$\pi$ の成分をエンドの射影写像と呼ぶ：
\[\pi_a \Colon \int_c p\ c\ c \to p\ a\ a\]
$\cat{C}$ が離散圏（恒等射以外の射がない）であれば、エンドは
圏 $\cat{C}$ 全体にわたる $p$ の全対角要素のグローバルな積に過ぎない。
後で、より一般的な場合においてエンドとこの積の間に等化子を通じた関係があることを示す。

Haskell では、エンドの式は全称量化子に直接変換される：

\src{snippet05}
厳密には、これは $p$ の全対角要素の積に過ぎないが、くさび条件は
\urlref{https://bartoszmilewski.com/2017/04/11/profunctor-parametricity/}{パラメトリシティ}により
自動的に満たされる。任意の関数 $f \Colon a \to b$ に対して、くさび条件は次のようになる：

\src{snippet06}
あるいは、型注釈付きで：

\begin{snipv}
dimap f id\textsubscript{b} . pi\textsubscript{b} = dimap id\textsubscript{a} f . pi\textsubscript{a}
\end{snipv}
ここで両辺の型は：

\src{snippet07}
\code{pi} は多相射影だ：

\src{snippet08}
ここでは、型推論が \code{e} の適切な成分を自動的に選択する。

コーンに対して可換条件の全集合を一つの自然変換として表現できたのと同様に、
くさび条件のすべてを一つの二重自然変換にまとめることができる。
そのためには、定数関手 $\Delta_c$ を、全ての対象の対を単一の対象 $c$ に写し、
全ての射の対をその対象の恒等射に写す定数プロ関手に一般化する必要がある。
くさびはその関手からプロ関手 $p$ への二重自然変換だ。
実際、$\Delta_c$ が全ての射を一つの恒等関数にリフトすることを認識すれば、
二重自然性の六角形はくさびのひし形に縮小する。

エンドは $\Set$ 以外の目標圏に対しても定義できるが、
ここでは $\Set$ 値プロ関手とそのエンドのみを考える。

\section{等化子としてのエンド}

エンドの定義における可換条件は等化子を使って書くことができる。
まず、二つの関数を定義しよう（この場合、数学的記法よりも使いやすいため
Haskell 記法を使う）。これらの関数はくさび条件の二つの収束する枝に対応する：

\src{snippet09}[b]
両関数はプロ関手 \code{p} の対角要素を以下の型の多相関数に写す：

\src{snippet10}
これらの関数は異なる型を持つ。しかし、\code{p} の全対角要素を集めた
一つの大きな積型を形成すれば、型を統一できる：

\src{snippet11}
関数 \code{lambda} と \code{rho} はこの積型からの二つの写像を導く：

\src{snippet12}
\code{p} のエンドはこれらの二つの関数の等化子だ。等化子は、
二つの関数が等しい最大の部分集合を選ぶことを思い出そう。
この場合、くさびの図式が可換になる全対角要素の積の部分集合を選ぶ。

\section{エンドとしての自然変換}

エンドの最も重要な例は、自然変換の集合だ。
二つの関手 $F$ と $G$ の間の自然変換は、
$\cat{C}(F a, G a)$ という形のホム集合から選ばれた射の族だ。
自然性条件がなければ、自然変換の集合はこれらのホム集合の積に過ぎない。
実際、Haskell では：

\src{snippet13}
Haskell でこれが機能する理由は、自然性がパラメトリシティから従うためだ。
しかし Haskell の外では、そのようなホム集合上の全ての対角的な切断が
自然変換をもたらすわけではない。しかし、写像：
\[\langle a, b \rangle \to \cat{C}(F a, G b)\]
がプロ関手であることに注意すれば、そのエンドを研究することは意義深い。
これがくさび条件だ：

\begin{figure}[H]
  \centering
  \includegraphics[width=0.4\textwidth]{images/end1.jpg}
\end{figure}

\noindent
集合 $\int_c \cat{C}(F c, G c)$ から一つの要素を選んでみよう。
二つの射影はこの要素を特定の変換の二つの成分に写す。それらを：
\begin{align*}
  \tau_a & \Colon F a \to G a \\
  \tau_b & \Colon F b \to G b
\end{align*}
と呼ぼう。左の枝では、ホム関手を使って射の対
$\langle \idarrow[a], G f \rangle$ をリフトする。
このようなリフトは同時前後合成として実装されることを思い出そう。
$\tau_a$ に作用すると、リフトされた対は：
\[G f \circ \tau_a \circ \idarrow[a]\]
を与える。図式の他の枝は：
\[\idarrow[b] \circ \tau_b \circ F f\]
を与える。くさび条件が要求するこれらの等式は、$\tau$ の自然性条件に他ならない。

\section{コエンド}

期待通り、エンドの双対をコエンドと呼ぶ。
これは\newterm{余くさび}（cowedge）と呼ばれるくさびの双対から構成される
（cow-edge ではなく co-wedge と発音する）。

\begin{figure}[H]
  \centering
  \includegraphics[width=0.25\textwidth]{images/end-31.jpg}
  \caption{An edgy cow?}
\end{figure}

\noindent
コエンドの記号は積分記号で、「積分変数」が上付きの位置に来る：
\[\int^c p\ c\ c\]
エンドが積に関係するように、コエンドは余積、すなわち和に関係する
（この点で、和の極限である積分に類似する）。
射影を持つのではなく、プロ関手の対角要素からコエンドへの入射がある。
余くさびの条件がなければ、プロ関手 $p$ のコエンドは $p\ a\ a$ か
$p\ b\ b$ か $p\ c\ c$ などどれかだと言えるだろう。
あるいは、コエンドがちょうど集合 $p\ a\ a$ となるような $a$ が存在すると言えるだろう。
エンドの定義で使った全称量化子は、コエンドでは存在量化子に変わる。

このため、疑似 Haskell ではコエンドを次のように定義する：

\begin{snip}{text}
exists a. p a a
\end{snip}
Haskell で存在量化子をエンコードする標準的な方法は、
全称量化されたデータコンストラクタを使うことだ。したがって次のように定義できる：

\src{snippet14}
この背景にある論理は、どのような $a$ を選んでも、
型の族 $p\ a\ a$ の値を使ってコエンドを構成できるはずだということだ。

エンドが等化子を使って定義できるように、コエンドは
\newterm{余等化子}（coequalizer）を使って記述できる。
全ての余くさびの条件は、可能な全ての関数 $b \to a$ に対して
$p\ a\ b$ の巨大な余積をとることでまとめられる。
Haskell では、これは存在型として表現される：

\src{snippet15}
この和型を評価する方法が二つあり、\code{dimap} を使って関数をリフトして
プロ関手 $p$ に適用する：

\src{snippet16}
ここで \code{DiagSum} は $p$ の対角要素の和だ：

\src{snippet17}
これら二つの関数の余等化子がコエンドだ。余等化子は、
\code{lambda} または \code{rho} を同じ引数に適用して得られる値を同一視することで
\code{DiagSum p} から得られる。ここで、引数は関数 $b \to a$ と
$p\ a\ b$ の要素からなる対だ。\code{lambda} と \code{rho} の適用は
型 \code{DiagSum p} の二つの潜在的に異なる値を生成する。
コエンドでは、この二つの値が同一視され、余くさびの条件が自動的に満たされる。

集合内の関連する要素を同一視するプロセスは、形式的には
商をとると呼ばれる。商を定義するには、
\newterm{同値関係}（equivalence relation）$\sim$——反射的、対称的、推移的な関係——が必要だ：
\begin{align*}
   & a \sim a                                                         \\
   & \text{if}\ a \sim b\ \text{then}\ b \sim a                       \\
   & \text{if}\ a \sim b\ \text{and}\ b \sim c\ \text{then}\ a \sim c
\end{align*}
このような関係は集合を同値類に分割する。各類は互いに関係する要素からなる。
各類から一つの代表元を選ぶことで\newterm{商集合}（quotient set）を形成する。
古典的な例は、整数の対（第二は非零）として有理数を定義することで、
以下の同値関係を持つ：
\[(a, b) \sim (c, d)\ \text{iff}\ a * d = b * c\]
これが同値関係であることは簡単に確認できる。
対 $(a, b)$ は分数 $\frac{a}{b}$ として解釈され、
分子と分母に公約数を持つ分数は同一視される。
有理数はそのような分数の同値類だ。

極限と余極限の以前の議論から、ホム関手が連続的であること、
すなわち極限を保存することを思い出すかもしれない。
双対として、反変ホム関手は余極限を極限に変換する。
これらの性質は、極限と余極限をそれぞれ一般化したエンドとコエンドにも一般化できる。
特に、コエンドをエンドに変換するための非常に有用な等式が得られる：
\[\Set(\int^x p\ x\ x, c) \cong \int_x \Set(p\ x\ x, c)\]
疑似 Haskell で見てみよう：

\begin{snipv}
(exists x. p x x) -> c \ensuremath{\cong} forall x. p x x -> c
\end{snipv}
これは、存在型をとる関数が多相関数と同値であることを示している。
このことは完全に理にかなっている。そのような関数は、
存在型にエンコードされている可能性のある型のいずれかを処理できる必要があるためだ。
これは、和型を受け取る関数が、和の中に存在する全ての型に対して一つずつハンドラのタプルを持つ
case 節として実装しなければならないという原則と同じだ。
ここでは、和型がコエンドに、ハンドラの族がエンドすなわち多相関数になる。

\section{忍者の米田の補題}

米田の補題に現れる自然変換の集合はエンドを使ってエンコードでき、
次の公式が得られる：
\[\int_z \Set(\cat{C}(a, z), F z) \cong F a\]
双対の公式もある：
\[\int^z \cat{C}(z, a)\times{}F z \cong F a\]
この等式は、ディラックのデルタ関数（関数 $\delta(a - z)$、あるいはより正確には
$a = z$ で無限の峰を持つ分布）の公式を強く連想させる。
ここでは、ホム関手がデルタ関数の役割を果たす。

これら二つの等式をまとめて\newterm{忍者の米田の補題}（Ninja Yoneda lemma）
と呼ぶことがある。

二番目の公式を証明するために、二つの対象が同型であるための必要十分条件は
それらのホム関手が同型であるという、米田の埋め込みの系を使う。
言い換えれば、$a \cong b$ であるための必要十分条件は、型：
\[[\cat{C}, \Set](\cat{C}(a, -), \cat{C}(b, =))\]
の自然変換が同型変換であることだ。

まず、証明したい等式の左辺を任意の対象 $c$ に向かうホム関手の内部に挿入することから始める：
\[\Set(\int^z \cat{C}(z, a)\times{}F z, c)\]
連続性の議論を使えば、コエンドをエンドに置き換えることができる：
\[\int_z \Set(\cat{C}(z, a)\times{}F z, c)\]
積と指数の間の随伴を活用できる：
\[\int_z \Set(\cat{C}(z, a), c^{(F z)})\]
米田の補題を使って「積分を実行」できる：
\[c^{(F a)}\]
（ここでは、関手 $c^{(F z)}$ が $z$ に関して反変であるため、
反変版の米田の補題を使ったことに注意。）
この指数対象はホム集合と同型だ：
\[\Set(F a, c)\]
最後に、米田の埋め込みを活用して同型に到達する：
\[\int^z \cat{C}(z, a)\times{}F z \cong F a\]

\section{プロ関手の合成}

プロ関手が関係を記述するというアイデアをさらに探求しよう——
より正確には、証明関連関係として、集合 $p\ a\ b$ が $a$ が $b$ に
関係するという証明の集合を表す。
二つの関係 $p$ と $q$ があれば、それらを合成しようとできる。
$q$ の後に $p$ を合成した結果を通じて $a$ が $b$ に関係するとは、
$q\ b\ c$ と $p\ c\ a$ の両方が空でないような
中間対象 $c$ が存在することだと言う。
この新しい関係の証明は個々の関係の証明の全ての対だ。
したがって、存在量化子がコエンドに対応し、二つの集合の直積が「証明の対」に対応するという
理解のもと、次の公式を使ってプロ関手の合成を定義できる：
\[(q \circ p)\ a\ b = \int^c p\ c\ a\times{}q\ b\ c\]
以下は \code{Data.Profunctor.Composition} からの同等の Haskell の定義（一部改名）だ：

\src{snippet18}
これは一般化代数的データ型（\acronym{GADT}）構文を使用しており、
自由型変数（ここでは \code{c}）は自動的に存在量化される。
（カリー化されていない）データコンストラクタ \code{Procompose} はしたがって
次と同値だ：

\begin{snip}{text}
exists c. (q a c, p c b)
\end{snip}
このように定義された合成の単位元はホム関手だ——
これは忍者の米田の補題から直ちに従う。
したがって、プロ関手が射として機能する圏が存在するかどうかという問いを立てることは
理にかなっている。答えは肯定的だが、プロ関手合成の結合性と単位元の法則は
自然同型を除いてのみ成立するという但し書きがある。
法則が同型を除いて成立するそのような圏を\newterm{双圏}（bicategory）と呼ぶ
（これは $\cat{2}$-圏よりも一般的だ）。
したがって、対象が圏、射がプロ関手、射の間の射（いわゆる2セル）が自然変換である
双圏 $\cat{Prof}$ がある。実際にはさらに進めることもできる。
プロ関手の他に、圏の間の射として通常の関手も存在するからだ。
二種類の射を持つ圏を\newterm{二重圏}（double category）と呼ぶ。

プロ関手は Haskell の lens ライブラリと arrow ライブラリにおいて重要な役割を果たす。
